\chapter{Masa Depan: Quantum AI & AGI}
\label{chap:future}

Apa yang terjadi dalam 5-10 tahun ke depan?

\section{Quantum Computing \& AI}

Komputer klasik berpikir dalam bit (0 atau 1). Komputer kuantum berpikir dalam Qubit (bisa 0 dan 1 sekaligus).
Jika digabungkan dengan AI, proses training model yang butuh berbulan-bulan bisa selesai dalam hitungan jam.

\section{Menuju AGI (Artificial General Intelligence)}

Kapan AI akan secerdas manusia dalam segala hal?
Prediksi ahli bervariasi dari 2027 hingga 2050.

\textbf{Tanda-tanda awal AGI:}
\begin{itemize}
    \item **Multimodality:** Bisa melihat, mendengar, dan bicara sekaligus.
    \item **Reasoning:** Bisa memecahkan masalah matematika yang belum pernah dilihat sebelumnya.
    \item **Long-term Memory:** Mengingat interaksi berbulan-bulan lalu.
\end{itemize}

\section{Regulasi AI (EU AI Act)}

Dunia mulai mengatur AI. Uni Eropa membagi AI menjadi 4 kategori risiko:
\begin{enumerate}
    \item **Unacceptable Risk:** Social scoring, manipulasi subliminal (Dilarang total).
    \item **High Risk:** Infrastruktur kritis, rekrutmen, penegakan hukum (Regulasi ketat).
    \item **Limited Risk:** Chatbot (Wajib transparansi "Saya adalah AI").
    \item **Minimal Risk:** Spam filter, video game (Bebas).
\end{enumerate}

\section{Penutup: Peran Manusia}

AI adalah alat. Kitalah pilotnya. Masa depan bukan milik AI, tapi milik manusia yang bekerja sama dengan AI.

\begin{center}
    \textit{Teruslah belajar, teruslah bereksperimen.}
\end{center}
